% Proposed B.121. Owned fragment only; no assembled source/register mutation.
\subsection*{The actual neck survives the four final prefix relations}

\statement{B.121}{Proved}{The actual six-cone mirror neck is a retract of an original-parity index-two subgroup after all four final relations; its free-five neck is a retract in the ordinary annular group}
\label{stmt:B121}
Keep the whole actual marking and intermediate group $M$ of B.120.
In particular $C$ is the adapted regional cone,
$c=aCa^{-1}$ is the selected gap cone, $K=F(a,B,d,e,f)$ with
$B=CbC$, and $P_i=O_1\cdots O_i$ are the actual five prefixes.
Let
\begin{equation}\label{eq:B121-actual-quotient}
 X=M/\langle\!\langle(jP_i)^2:\ 1\leq i\leq4\rangle\!\rangle,
 \qquad
 Q=K\rtimes_\eta C_2(j)
   =\mathop{*}_{i=0}^{5}C_2(r_i),
 \quad r_0=j,\ r_i=jP_i.
\end{equation}
The canonical homomorphism $\iota:Q\to X$ comes from the
actual neck and mirror, with every remaining relation imposed.
Let $\varepsilon:X\to C_2$ be the original ordinary-cover
character, extended by $\varepsilon(j)=0$, so
$\varepsilon(C)=1$ and it vanishes on the ordinary regional
generators. Let $\mu:X\to C_2$ be mirror parity, zero on $R$
and one on $j$. Write
\[
 X^0=\ker\varepsilon,\qquad
 \Gamma=\ker\mu,\qquad
 \Gamma^0=\ker\mu\cap\ker\varepsilon.
\]
Thus $\Gamma$ is the actual annular one-neck orbifold group and
$\Gamma^0$ is its actual ordinary double cover. The two
characters are distinct.

There is a literal retraction
\begin{equation}\label{eq:B121-retractions}
 \mathcal R:X^0\longrightarrow Q,
 \qquad \mathcal R\iota=1_Q,
 \qquad \mu|_Q\circ\mathcal R=\mu|_{X^0}.
\end{equation}
In particular $\iota$ is injective. Restricting this retraction
along mirror parity gives a retraction
\begin{equation}\label{eq:B121-ordinary-neck-retract}
 \mathcal R|_{\Gamma^0}:\Gamma^0\longrightarrow K,
 \qquad \mathcal R|_K=1_K.
\end{equation}
Consequently $Q$ has its full induced profinite topology in $X$
and $K$ has its full induced profinite topology in $\Gamma^0$
and $\Gamma$. The natural full completions therefore embed:
\begin{equation}\label{eq:B121-completed-embeddings}
 \widehat Q\hookrightarrow\widehat X,
 \qquad \widehat K\hookrightarrow\widehat{\Gamma^0}
                         \hookrightarrow\widehat\Gamma.
\end{equation}
This is a statement about the actual final quotient. It does not
assert that either regional vertex embeds after the final relations,
or that the source image is free, or that a selected-cone
centralizer is exhausted by the source.

\begin{proof}
\emph{B.121.1. The complete presentation and the two parities.}
B.120 gives the exact whole presentation
\begin{equation}\label{eq:B121-whole-presentation}
 X=\langle R,j\mid j^2=1,\ [jcj,k]=1,
                       \ jP_i j=P_i^{-1}\ (1\leq i\leq5)\rangle.
\end{equation}
The $j$-conjugate of the displayed cross commutator is the other
actual cross relation. Every regional node and retained-base
relation is in $R$; none is omitted. Each $P_i$ is ordinary,
so both $\varepsilon$ and $\mu$ respect this presentation with
the values stated above. The six-reflection presentation of $Q$
in \eqref{eq:B121-actual-quotient} is a reversible change from
the free prefix basis: $P_i=r_0r_i$, and
$jP_i j=P_i^{-1}$ is equivalent to $r_i^2=1$.
At this point $Q$ is an abstract group with its canonical map
to $X$; its injection will follow from the retraction rather
than being assumed.

\emph{B.121.2. The actual detector and all five conjugate prefixes.}
Use the actual ordinary regional retraction from B.120.7,
\[
 r:R^0\to K,\qquad r|_K=1_K,
\]
where $R^0=\ker(\varepsilon|_R)$. Its companion
$\zeta=r\circ\operatorname{Ad}(C)|_K$ is exactly
\begin{equation}\label{eq:B121-companion}
 \zeta(a)=\zeta(B)=\zeta(d)=\zeta(e)=1,
 \qquad \zeta(f)=f.
\end{equation}
This is the whole-region detector, already verified on all
ordinary node and both signed generic action relations. In
particular it is valid after every regional filling.
It kills $k=f^{-1}h$ and $CkC$, as proved in B.120.7.

The exact prefix expansions and their companion images are
\begin{equation}\label{eq:B121-prefix-images}
 \begin{array}{c|l|c}
 i&P_i&\zeta(P_i)=r(CP_iC)\\\hline
 1&aBa^{-1}&1\\
 2&aB&1\\
 3&aBfde d^{-1}f^{-1}&1\\
 4&aBfde d^{-1}f^{-1}d&1\\
 5&aBfde&f.
 \end{array}
\end{equation}
These are literal evaluations: for $i=3,4$ the two surviving
$f$ letters cancel, and no relation in the target free group
other than free cancellation is used. This simultaneous check
of all four additional prefixes is the extra input needed
after B.120's intermediate construction.

\emph{B.121.3. The exact ordinary presentation after killing the base loop.}
Put $\bar R=R/\langle\!\langle k\rangle\!\rangle_R$ and
kill the normal closure of $k$ in $X$. The resulting whole group
has the exact presentation
\begin{equation}\label{eq:B121-killed-X}
 \bar X=\langle\bar R,j\mid j^2=1,
                    \ jP_i j=P_i^{-1}\ (1\leq i\leq5)\rangle.
\end{equation}
The cross relation disappears because $k$ is killed. The original
character $\varepsilon$ is still onto, with representatives
$1,C$. The complete two-coset rewriting therefore gives
\begin{equation}\label{eq:B121-two-reflection-presentation}
 \begin{gathered}
 \bar X^0=\langle\bar R^0,j_0,j_1\mid j_0^2=j_1^2=1,
       \ j_0P_i j_0=P_i^{-1},\ j_1Q_i j_1=Q_i^{-1}
                                      \ (1\leq i\leq5)\rangle,\\
 j_0=j,\qquad j_1=CjC,\qquad Q_i=CP_iC.
 \end{gathered}
\end{equation}
There is no further relation between $j_0$ and $j_1$: these
are precisely the two rewritten copies of each added relation;
the old regional relations give $\bar R^0$.
Since $R=R^0\rtimes\langle C\rangle$, the normal closure of
$k$ in $R$ is the normal closure inside $R^0$ of $k$ and $CkC$.
The retraction $r$ kills both, so it factors as
$\bar r:\bar R^0\to K$, still fixing every element of $K$.
Its values on the $Q_i$ are exactly
\eqref{eq:B121-prefix-images}.

\emph{B.121.4. An actual reflection in the final neck inverting $f$.}
Within the abstract $Q$, put
\begin{equation}\label{eq:B121-inverting-reflection}
 t=d\,r_4\,d^{-1}=d(jP_4)d^{-1}.
\end{equation}
It is an involution. The actual basis relation
$O_5=d^{-1}fd=P_4^{-1}P_5$ gives
\[
 f=d(P_4^{-1}P_5)d^{-1}=d(r_4r_5)d^{-1}.
\]
Therefore, using only the two cone-square relations in $Q$,
\begin{equation}\label{eq:B121-inverts-f}
 tft=d(r_5r_4)d^{-1}=f^{-1}.
\end{equation}
This is an explicit word, not a conjugacy or existence assertion
about involutions in an unspecified quotient.

Define a map from \eqref{eq:B121-two-reflection-presentation}
to $Q$ by
\begin{equation}\label{eq:B121-explicit-map}
 \bar R^0\xrightarrow{\bar r}K\hookrightarrow Q,
 \qquad j_0\longmapsto j,\qquad j_1\longmapsto t.
\end{equation}
Every original regional relation is respected by $\bar r$.
The first five inversion relations hold in $Q$ by definition.
For $i<5$ the second five relations have $\bar r(Q_i)=1$
and are automatic; the last is precisely
\eqref{eq:B121-inverts-f}. The two involution relations also
hold. Thus \eqref{eq:B121-explicit-map} is a homomorphism on
the complete actual quotient, not a map defined only on a
generic fiber or a partial list of relators.

Compose with $X^0\to\bar X^0$ to obtain $\mathcal R$.
Its composite with the canonical $Q\to X^0$ fixes all of $K$
and $j$, hence is the identity on $Q$. This proves both the
injection of the actual mirror neck and the retraction in
\eqref{eq:B121-retractions}. Notice that this assigns $j_1$
to $t$ under a homomorphism; it does not claim that
$CjC=t$ is an identity in $X$.

\emph{B.121.5. Mirror compatibility and the actual ordinary neck.}
Every element in the image of the regional subgroup $\bar R^0$
has mirror parity zero, and its image under $\bar r$ lies in $K$,
also of parity zero. Both $j_0$ and $j_1=CjC$ have mirror
parity one. Their target images have parity one as well:
$\mu(j)=1$ and $t=d(jP_4)d^{-1}$ has parity one because
$d,P_4\in K$. Hence \eqref{eq:B121-explicit-map} preserves
mirror parity, proving the last identity in
\eqref{eq:B121-retractions}.

In the abstract semidirect product $Q=K\rtimes C_2(j)$, the
mirror-parity kernel is exactly $K$. Consequently the retraction
restricts to \eqref{eq:B121-ordinary-neck-retract}. The group
on its left is exactly the actual ordinary one-neck group:
first take the actual annular orbifold group $\Gamma=\ker\mu$,
then its original ordinary double cover. The argument does not
identify these two independent parity maps.

\emph{B.121.6. Full topologies and exact finite-kernel refinements.}
For every $V\lhd Q$ of finite index, the subgroup
$\mathcal R^{-1}(V)\lhd X^0$ has finite index and intersects
the embedded $Q$ exactly in $V$. Core it through the actual
index-two extension $X^0\leq X$. Since $C^2=1$, this core is
the intersection of that subgroup and its $C$-conjugate. It is
normal of finite index in $X$ and its intersection with $Q$ is
contained in $V$. This is a cofinal finite-kernel refinement
proving the full induced topology of $Q$ in $X$.

Likewise a prescribed finite-index normal subgroup of $K$ pulls
back under \eqref{eq:B121-ordinary-neck-retract} to a finite-index
normal subgroup of $\Gamma^0$ with exactly that intersection.
Taking its two-coset core in $\Gamma$ gives the full topology
there too. The maps of full completions in
\eqref{eq:B121-completed-embeddings} are therefore injective.
The retractions themselves extend continuously to the indicated
full completions. In particular the embedded completed neck is
the full profinite free group $\widehat{F_5}$, not only a
pro-$p$ or prosolvable image.

Finally, this retraction is not claimed to commute with the
original selected-cone action. The image of a fixed element
under it need not identify a source preimage. Thus the actual
neck embedding and topology proved here leave the full selected
source-centralizer comparison, and the effect of the final
relations on the other regional generators, unresolved.
\end{proof}
